% ============================================================================
%  3.1 ANÁLISIS DE LA NECESIDAD
% ============================================================================

\section{An\'alisis de la necesidad}

\subsection{Contexto del proyecto}

La aplicaci\'on \nombreApp{} es un sistema m\'ovil para la
visualizaci\'on interactiva de escenarios deportivos a nivel nacional
en Bolivia. Permite a los ciudadanos localizar y acceder a informaci\'on
confiable sobre infraestructura deportiva del pa\'is, incluyendo
estadios, coliseos, canchas m\'ultiples y polideportivos.

El sistema actual cuenta con las siguientes funcionalidades principales:

\begin{itemize}
  \item \textbf{Mapa interactivo:} Ubicaci\'on de escenarios en un mapa
    con MapLibre (nativo y web).
  \item \textbf{Detalle de escenario:} Informaci\'on completa incluyendo
    imagenes, deportes disponibles, horarios y eventos pr\'oximos.
  \item \textbf{Favoritos:} Guardado de escenarios de inter\'es del
    usuario.
  \item \textbf{B\'usqueda:} B\'usqueda por nombre, deporte o
    ubicaci\'on.
  \item \textbf{Gesti\'on:} CRUD de escenarios para roles admin y gestor.
  \item \textbf{Visor 360:} Exploraci\'on panor\'amica de sectores del
    estadio.
  \item \textbf{Eventos:} Creaci\'on y consulta de eventos deportivos
    asociados a escenarios.
\end{itemize}

\subsection{El problema identificado}

Si bien la aplicaci\'on permite \textit{ver} los eventos pr\'oximos
asociados a cada escenario, no existe un mecanismo para que los
usuarios \textbf{se inscriban} a dichos eventos. Esto genera las
siguientes limitaciones:

\begin{enumerate}
  \item \textbf{Sin seguimiento:} El usuario no tiene forma de indicar
    qu\'e eventos le interesan asistir.
  \item \textbf{Sin conteo de asistencia:} No se conoce cu\'antas
    personas planean asistir a cada evento.
  \item \textbf{Sin recordatorios:} No hay base para futuras
    notificaciones de eventos pr\'oximos.
  \item \textbf{Experiencia incompleta:} La app muestra eventos pero
    no permite interactuar con ellos de manera significativa.
\end{enumerate}

\subsection{Justificaci\'on del nuevo m\'odulo}

La incorporaci\'on de un m\'odulo de inscripci\'on a eventos resuelve
estas limitaciones y aporta valor tanto al usuario final como al
gestor del escenario:

\begin{itemize}
  \item \textbf{Para el usuario:} Puede inscribirse a eventos de
    inter\'es, cancelar su inscripci\'on, y consultar sus inscripciones
    desde su perfil.
  \item \textbf{Para el gestor:} Puede ver cu\'antas personas est\'an
    inscritas en cada evento, lo que le permite planificar
    log\'istica y capacidad.
  \item \textbf{Para el admin:} Tiene visibilidad completa de todas las
    inscripciones del sistema.
\end{itemize}
